\item \label{itm:hpmin} {\bf [5 pts.] Lord Kelvin and the age of the Earth?} \\
In the late 1800's William Thompson (Lord Kelvin) estimated the age of the Earth
between 24 and 400 million years.  He did so due to assumptions about the
cooling rate of the Earth.

\begin{enumerate}
    \item If the Earth were infinitely old, what would the temperature gradient
    of the interior be?
    
    \item The thermal gradient of the continents are typically between 10 and
    30$^\circ$C~km$^{-1}$.  Using these thermal gradients as bounds, estimate
    the age of the Earth assuming the initial temperature of the Earth is
    approximately that of completely molten rock ($\sim$3900$^\circ$C).

    \item If instead, we assume the Earth has been cooling for 4.5~Ga, how deep
    will cooling progress?

    \item What does Kelvin's logic fail?  What are additional possible physical
    influnces on cooling that are not captured by Kelvin's simple analysis that
    would have a significant on cooling?
\end{enumerate}
\FloatBarrier

